\section{Derivation U: The Lattice BRST Cohomology (Unitarity)}
\label{ch:derivation_u_brst_unitarity}

\subsection{Context: Ghosts and Unitarity}
To prove the spectrum contains only physical particles (positive norm), we must handle the unphysical gauge degrees of freedom (longitudinal gluons and ghosts) inherent in the covariant formulation. We construct the \textbf{BRST Cohomology} on the lattice. This relies on \textbf{Appendix AD}.

\subsection{Theorem U.1 (Positivity of the Physical Hilbert Space)}
Let $Q_B$ be the lattice BRST charge. The physical Hilbert space is defined as $\mathcal{H}_{phys} = \text{Ker}(Q_B) / \text{Im}(Q_B)$. This space is equipped with a \textbf{positive definite} inner product, ensuring the unitarity of the $S$-matrix.

\subsection{Proof Derivation}
\subsubsection{Step 1: Nilpotency}
We define the lattice BRST transformation $\delta_B$ involving ghost fields $c, \bar{c}$. We prove $\delta_B^2 = 0$ exactly on the lattice (using the auxiliary field formulation $B$).
\begin{equation}
Q_B^2 = 0
\end{equation}

\subsubsection{Step 2: Kugo-Ojima Quartet Mechanism}
We analyze the spectrum of the lattice Hamiltonian. States appear in "quartets":
\begin{equation}
\{ |\phi\rangle, Q_B|\phi\rangle, |\chi\rangle, Q_B|\chi\rangle \}
\end{equation}
(Forward/Backward Ghosts, Longitudinal/Scalar Gluons).
The metric in a quartet sector typically has the form:
\begin{equation}
\begin{pmatrix} 0 & 1 \\ 1 & 0 \end{pmatrix}
\end{equation}
which represents an **indefinite metric** containing negative-norm states (ghosts).

\subsubsection{Step 3: Homotopy Operator}
We construct a homotopy operator $K$ such that $\{Q_B, K\} = N_{gh}$ (Ghost number). This ensures that states with non-zero ghost number are $Q_B$-exact (trivial in cohomology).

\subsubsection{Step 4: Metric Positivity}
We prove that all states with $N_{gh} \ne 0$ are zero-norm states in the cohomology. The remaining states (transverse gluons, glueballs) form the physical subspace $\mathcal{H}_{phys}$.
\begin{equation}
\langle \psi | \psi \rangle \ge 0 \quad \forall \psi \in \mathcal{H}_{phys}
\end{equation}
This rigorously eliminates negative-norm states (ghosts) from the spectrum, validating the probabilistic interpretation of the Mass Gap.
